\section{Which coefficients remain free}\label{app:freedom}

Throughout this appendix assume \(a^{p-1}=0\). The finite formula for \(L_a\) gives one working coordinate. We can also find every other coordinate beginning with \(t\).

First solve a simpler coefficient problem: which expressions \(P(t)\), with no constant term, satisfy
\[
 P(x+y)=P(x)+P(y)
\]
in arithmetic modulo \(p\)? Write \(P(t)=\sum_{n\ge1}q_nt^n\). Comparing a mixed coefficient gives
\begin{equation}\label{eq:additive-coefficients}
 q_n\binom ni=0\qquad(0<i<n).
\end{equation}
If any interior binomial coefficient is nonzero modulo \(p\), it has a modular inverse and forces \(q_n=0\).

The indices with no such mixed coefficient are precisely powers of \(p\). Here is a direct check. Write the base-\(p\) digits of \(n\) as
\[
 n=n_0+n_1p+n_2p^2+\cdots,\qquad0\le n_j<p.
\]
Repeatedly using \((u+v)^p=u^p+v^p\) gives the polynomial identity
\[
 (1+z)^n=\prod_j(1+z^{p^j})^{n_j}
\]
modulo \(p\). Let \(r\) be the least index with \(n_r\ne0\). The coefficient of \(z^{p^r}\) on the right is \(n_r\ne0\): it comes from choosing the linear term of \((1+z^{p^r})^{n_r}\). All other nonconstant choices have larger degree. Unless \(n=p^r\), this is an interior term, so \eqref{eq:additive-coefficients} forces \(q_n=0\). When \(n=p^r\), repeated cubing or, more generally, taking \(p\)-th powers shows that no mixed terms occur.

Consequently all the expressions that preserve addition are exactly
\[
 P(t)=b_0t+b_1t^p+b_2t^{p^2}+\cdots.
\]
They begin with \(t\) when \(b_0=1\).

Now let \(M\) be any coordinate beginning with \(t\) that makes \(F_a\) additive. Use the explicit inverse \(E_a\) and put \(P(s)=M(E_a(s))\). The inverse identities proved earlier give
\[
 E_a(x+y)=F_a(E_a(x),E_a(y)).
\]
Applying \(M\) to this equality proves \(P(x+y)=P(x)+P(y)\). Therefore every possible \(M\) is obtained by following \(L_a\) with one of the addition-preserving expressions above.

The resulting formula simplifies. In each nonlinear term of \(L_a(t)\), the coefficient contains \(a^{n-1}\) with \(n\ge2\). Raising it to the \(p^r\)-th power gives a factor \(a^{(n-1)p^r}=0\), because \((n-1)p^r\ge p-1\). Thus
\[
 L_a(t)^{p^r}=t^{p^r}\qquad(r\ge1),
\]
and all the working coordinates beginning with \(t\) are
\begin{equation}\label{eq:all-logarithms}
 \boxed{\;M(t)=L_a(t)+b_1t^p+b_2t^{p^2}+\cdots,\;}
\end{equation}
with the \(b_r\) freely chosen in the coefficient arithmetic. Conversely, each expression in \eqref{eq:all-logarithms} works. For the added powers, expand directly:
\[
 F_a(x,y)^{p^r}
 =x^{p^r}+y^{p^r}+a^{p^r}x^{p^r}y^{p^r}
 =x^{p^r}+y^{p^r}.
\]
The last term vanishes because \(a^{p-1}=0\) and \(p^r\ge p-1\). Each added term therefore has zero error for this particular operation.

An expression beginning with \(t\) always has an inverse as a series, by the recursive coefficient construction in Section~\ref{sec:coordinate-action}. That inverse need not stop after finitely many terms. For example, \(t+t^p\) over ordinary numbers modulo \(p\) cannot have a polynomial inverse: composing it with a nonconstant polynomial of degree \(d\) gives degree \(pd\), which cannot be one. The particular pair \(L_a,E_a\) has the stronger property that both expressions are finite polynomials.
